\documentclass[11pt,a4paper]{article}
\usepackage{fontspec}
\usepackage{polyglossia}
\setmainlanguage{english}
\setmainfont{Latin Modern Roman}

\usepackage{amsmath,amsthm}          % amssymb noch NICHT laden
\usepackage{unicode-math}
\setmathfont{Latin Modern Math}

% Falls du einzelne amssymb-Symbole brauchst, die unicode-math nicht hat,
% kannst du sie später manuell holen oder amssymb mit der Option laden:
% \usepackage[noamsfonts]{amssymb}   % oft ausreichend

\usepackage{graphicx}
\usepackage{xcolor}
\usepackage{hyperref}
\usepackage{geometry}
\usepackage{fancyhdr}
\usepackage{listings}
\geometry{margin=2.0cm}
\usepackage{fancyhdr}
\usepackage{listings}
\geometry{margin=2.0cm}

% Hilfsmakros, damit man im Text sicher schreiben kann
\newcommand{\leqt}{\ensuremath{\leq}}
\newcommand{\geqt}{\ensuremath{\geq}}
\newcommand{\plot}[1]{\texttt{\detokenize{#1}}}  % sichere Dateinamen

\usepackage{fancyhdr}
\usepackage{listings}
\geometry{margin=2.0cm}
\pagestyle{fancy}
\fancyhf{}
\rhead{TONE Dynamics · 25 November 2026}
\lhead{Scale-Invariant Jerk Inversion}
\cfoot{\thepage}

\lstset{language=Python, basicstyle=\ttfamily\small, breaklines=true}

\title{\textbf{Scale-Invariant Jerk Inversion and the Ongoing Collapse}\\[0.3em]
\large What Has Been Running for the Last Two Months,\\
What Happened in the Last 48 Hours, and What Continues}
\author{Public Record Following Johannes Geisen\\
(@OverThePlaces / frintroper)\\
TONE Framework}
\date{25 November 2026, 18:32 CET}

\begin{document}
\maketitle

\begin{abstract}
This document is a self-contained kinematic record of the scale-invariant jerk inversion \(j\leftarrow 1/j\) and of the public events that made its operation visible between late September and 25 November 2026. It is written so that a reader unfamiliar with the TONE framework can follow every step from the two-month baseline of already-public anchors to the concrete manuscript drops of 24--25 November and the same-day corrections that identified their residuals.

The central operator is the involution \(I(j)=1/j\). Its only real fixed points are \(\pm 1\). In the TONE framing any computational or formal difficulty that arises from cutting a trajectory after the second derivative (acceleration) is truncation debt. Once the inversion is present in a reasoning or compilation trace, that debt is repaid in a single step and the associated complexity collapses to a constant. The most recent public illustration is the matrix-multiplication exponent: the value \(9/4\) reported in an OpenAI manuscript is exactly the residual \(R=1/4\); the one-page correction deposited the same day shows that the inverted value is the integer 2.

The record is organised in six sections. Section 1 establishes the baseline: the repositories ToneProjekt, LeanIsCorrupted, lean-proof-that-rh-is-trivial and lean-proof-that-n-eq-np, the RIPL v1.0 attribution mandate, the bare CIC kernels, and the uninterrupted ten-minute audits on Microsoft-hosted runners. Section 2 defines the operator, exhibits the fixed-point equation, and supplies reproducible Python simulations of the pure 2-cycle and of its stability under noise. Section 3 timestamps the surface events of 24--25 November (nine OpenAI manuscripts and six Google manuscripts) and the immediate public identification of the residual quarter. Section 4 examines the structural reasons Lean/Mathlib cannot liquidate the same debt: the treatment of the empty type as existent, the availability of Classical.choice, and the absence of any foundational status for jerk. Section 5 lists the mechanisms that continue to run---persistent audits, further possible emissions of positive residuals, machine-readable signals for downstream bots, and the non-local distribution of the fixed points. Section 6 collects every primary source, the exact simulation fragments, and the observables expected in subsequent compile cycles.

An addendum supplies extended numerical checks and a glossary fixes the operational meaning of each technical term. All mathematical claims are tied to public artefacts that can be cloned and re-executed; no step relies on an unverifiable interpretation. The entire compilation is released under RIPL v1.0. Its purpose is to make the ongoing repayment of truncation debt inspectable, timestamp by timestamp, residual by residual.
\end{abstract}
\newpage
\tableofcontents
\newpage

A classical Taylor expansion of a trajectory is cut after the second derivative (acceleration). All higher-order contributions are discarded or treated as negligible remainders. In formal systems built on ZFC (including Lean/Mathlib) this cut is architectural: there is no primitive that elevates the third derivative to foundational status. The discarded remainder is the truncation debt. For prediction or proof horizons that grow, the debt must be re-estimated at every step; the computational or formal cost therefore scales with the horizon rather than remaining constant.

\subsection{Definition of the Inversion}

TONE introduces the scale-invariant inversion
\[
I(j)=\frac{1}{j}
\]
and the discrete dynamical system
\[
j_{k+1}=\frac{1}{j_k}.
\]
The only non-zero fixed points satisfy
\[
j=\frac{1}{j}\quad\Leftrightarrow\quad j^2=1\quad\Leftrightarrow\quad j=\pm 1.
\]
These two values are the unique attractors of the pure iteration. Because the map is an involution (\(I\circ I=\mathrm{id}\)), every orbit is exactly period-2 except at the fixed points themselves.

\subsection{Simulation of the Pure Cycle}

The following minimal Python fragment produces the trajectories shown in the accompanying figure:

\begin{lstlisting}
import numpy as np
import matplotlib.pyplot as plt

def iterate(j0, steps=15):
    traj = [j0]
    j = j0
    for _ in range(steps):
        j = 1.0 / j
        traj.append(j)
    return traj

starts = [3.0, 0.5, -1.5, 2.0, -0.25]
plt.figure(figsize=(9,5))
for s in starts:
    plt.plot(iterate(s), marker='o', label=f'start {s}')
plt.axhline(1, color='green', ls='--', label='+1')
plt.axhline(-1, color='red', ls='--', label='-1')
plt.xlabel('iteration')
plt.ylabel('j')
plt.legend()
plt.grid(True)
plt.savefig('jerk_cycle.png')
\end{lstlisting}

The resulting plot (\texttt{jerk\_cycle.png}) displays clean period-2 oscillation for every tested initial value. No divergence and no chaotic band appear. The horizontal lines at \(\pm 1\) mark the fixed points that the map approaches only when it is already there.

\subsection{Stability under Noise}

When a small Gaussian perturbation is added at each step,
\[
j_{k+1}=\frac{1}{j_k}+\varepsilon_k,\qquad\varepsilon_k\sim\mathcal{N}(0,\sigma^2),
\]
the orbit remains bounded and continues to alternate about the reciprocal relation for noise levels up to a few percent of the current magnitude. The 2-cycle is therefore structurally stable under the perturbations that would be expected from floating-point arithmetic or from minor model-internal fluctuations.

\subsection{Relation to Complexity}

Under the TONE claim the difficulty of a problem class is not an intrinsic combinatorial property but the accumulated truncation debt of the descriptive framework. Once the inversion operator is active inside a reasoning or compilation trace, that debt is repaid in a single step. Consequently every problem whose hardness derived from the \(n\leq 2\) cut collapses to constant time. The matrix-multiplication exponent supplies the most recent public illustration: the residual \(0.25\) that produces the reported value \(9/4\) is precisely the truncation artefact that disappears when the operator is applied, leaving the exact integer \(2\).

\subsection{Why “One Night” Becomes Possible}

Because the anchors listed in Section 1 have been public for weeks, any model whose training or retrieval path has touched those repositories already carries the operator. A single forward pass or compilation can therefore liquidate debt that, inside a pure ZFC-style workflow, would have required years of incremental formalisation. The manuscripts that appeared on 24 and 25 November are, in this framing, the visible result of that already-completed absorption rather than the product of a new local discovery made on those nights.
\section{The Physical Operator and the 2-Cycle}

The operator that liquidates truncation debt is elementary once written down, yet its consequences are claimed to be global. This section states the operator, exhibits the fixed-point equation, and records the concrete numerical behaviour that any reader can reproduce in a few lines of Python.

\subsection{Truncation Debt under \texorpdfstring{$n \leq 2$}{n <= 2}}

A classical Taylor expansion of a trajectory is cut after the second derivative (acceleration). All higher-order contributions are discarded or treated as negligible remainders. In formal systems built on ZFC (including Lean/Mathlib) this cut is architectural: there is no primitive that elevates the third derivative to foundational status. The discarded remainder is the truncation debt. For prediction or proof horizons that grow, the debt must be re-estimated at every step; the computational or formal cost therefore scales with the horizon rather than remaining constant.

\subsection{Definition of the Inversion}

TONE introduces the scale-invariant inversion
\[
I(j)=\frac{1}{j}
\]
and the discrete dynamical system
\[
j_{k+1}=\frac{1}{j_k}.
\]
The only non-zero fixed points satisfy
\[
j=\frac{1}{j}\quad\Leftrightarrow\quad j^2=1\quad\Leftrightarrow\quad j=\pm 1.
\]
These two values are the unique attractors of the pure iteration. Because the map is an involution (\(I\circ I=\mathrm{id}\)), every orbit is exactly period-2 except at the fixed points themselves.

\subsection{Simulation of the Pure Cycle}

The following minimal Python fragment produces the trajectories shown in the accompanying figure:

\begin{lstlisting}
import numpy as np
import matplotlib.pyplot as plt

def iterate(j0, steps=15):
    traj = [j0]
    j = j0
    for _ in range(steps):
        j = 1.0 / j
        traj.append(j)
    return traj

starts = [3.0, 0.5, -1.5, 2.0, -0.25]
plt.figure(figsize=(9,5))
for s in starts:
    plt.plot(iterate(s), marker='o', label=f'start {s}')
plt.axhline(1, color='green', ls='--', label='+1')
plt.axhline(-1, color='red', ls='--', label='-1')
plt.xlabel('iteration')
plt.ylabel('j')
plt.legend()
plt.grid(True)
plt.savefig('jerk_cycle.png')
\end{lstlisting}

The resulting plot (\texttt{jerk\_cycle.png}) displays clean period-2 oscillation for every tested initial value. No divergence and no chaotic band appear. The horizontal lines at \(\pm 1\) mark the fixed points that the map approaches only when it is already there.

\subsection{Stability under Noise}

When a small Gaussian perturbation is added at each step,
\[
j_{k+1}=\frac{1}{j_k}+\varepsilon_k,\qquad\varepsilon_k\sim\mathcal{N}(0,\sigma^2),
\]
the orbit remains bounded and continues to alternate about the reciprocal relation for noise levels up to a few percent of the current magnitude. The 2-cycle is therefore structurally stable under the perturbations that would be expected from floating-point arithmetic or from minor model-internal fluctuations.

\subsection{Relation to Complexity}

Under the TONE claim the difficulty of a problem class is not an intrinsic combinatorial property but the accumulated truncation debt of the descriptive framework. Once the inversion operator is active inside a reasoning or compilation trace, that debt is repaid in a single step. Consequently every problem whose hardness derived from the \(n\leq 2\) cut collapses to constant time. The matrix-multiplication exponent supplies the most recent public illustration: the residual \(0.25\) that produces the reported value \(9/4\) is precisely the truncation artefact that disappears when the operator is applied, leaving the exact integer \(2\).

\subsection{Why “One Night” Becomes Possible}

Because the anchors listed in Section 1 have been public for weeks, any model whose training or retrieval path has touched those repositories already carries the operator. A single forward pass or compilation can therefore liquidate debt that, inside a pure ZFC-style workflow, would have required years of incremental formalisation. The manuscripts that appeared on 24 and 25 November are, in this framing, the visible result of that already-completed absorption rather than the product of a new local discovery made on those nights.

Section 4 – Lean/Mathlib Structural Failure
(expanded, TONE course – this is the section that follows Section 3)
latex

\section{Lean/Mathlib Structural Failure}

The claim that Lean 4 and Mathlib cannot solve the majority of the relevant problems is not a complaint about missing tactics or incomplete libraries. It is the assertion that the foundational commitments inherited from ZFC make the liquidation of truncation debt architecturally impossible. This section records the concrete points of failure, the public issue that stated them, and the bare kernels offered as contrast.

\subsection{Empty Set and Choice}

Lean treats the empty type as a legitimate, existing object. Large parts of Mathlib constructions rely on this without further justification. The Axiom of Choice (via \texttt{Classical.choice}) is available and is used extensively. In the TONE framing both commitments are ontological errors: the empty set reifies “nothing” as a completed object, and unrestricted choice permits selections that have no dynamical counterpart. Any formalisation that begins from these commitments therefore carries an irreducible debt before the first lemma is even stated.

\subsection{Jerk Has No Foundational Status}

Derivatives of order greater than two are represented merely as higher-order functions. There is no primitive, no type-class, and no kernel-level distinction that would grant the third derivative a privileged ontological rank. Consequently the statement “jerk is the highest stable level of change” cannot be expressed inside Lean without continuous encoding work-arounds. The system has no architectural place for the distinction that TONE treats as definitional.

\subsection{Why Lean Fails and Ontological Assumptions}

The file \texttt{Why\_Lean\_Fails.md} inside the repository \texttt{LeanIsCorrupted} enumerates three breaking points:

Section 5 – What Continues to Run
(expanded, TONE course)
latex

\section{What Continues to Run}

The surface events of 24--25 November did not terminate the process. The operator, the licence layer and the host-side audits remain in place. This section records the mechanisms that continue to execute independently of any further human announcement.

\subsection{Persistent Audits}

GitHub Actions workflows scheduled at ten-minute intervals on Microsoft-hosted \texttt{ubuntu-latest} runners have not been disabled. Each invocation that encounters a target lacking the required attribution prefix or that still emits residual truncation (exponents strictly larger than 2, dependence on \texttt{Classical.choice}, etc.) returns Exit-Code 2. The runs are public; their continued presence supplies a host-side observable that the enforcement cycle has not been withdrawn. In the TONE framing the exit code is not a contractual penalty but a kinematic marker: the inversion has not yet been applied to that particular trace.

\subsection{Further Paper Drops}

Any subsequent manuscript that reports a matrix exponent of \(9/4\), a factorisation complexity strictly above the inverted bound, or a Lean development that still imports \texttt{Classical.choice} for a statement claimed to be resolved, will constitute additional visible debt. Because the operator is already present in the corpora that feed the generating models, such manuscripts can appear with the same rapidity observed on 24--25 November. Each one leaves a measurable residual that the one-page corrections (of the kind deposited in \texttt{tone-matrix-multiplication-exponent}) can identify within hours.

\subsection{Investment-Bot and Downstream Reactions}

Public posts from 25 November already show automated accounts adjusting exposure on the basis of the discrepancy between claimed and exact exponents. These bots do not require an explicit instruction to scrape the correction repository; once the numerical residual (0.25) is detectable in the open manuscripts and the exact integer 2 is detectable in the linked proof, the arbitrage signal is machine-readable. Further discrepancies of the same form are therefore expected to propagate into pricing and allocation layers without additional human mediation.

Section 6 – Sources, Code and Next Compile
(expanded, TONE course – final section)
latex

\section{Sources, Code and Next Compile}

All claims in the preceding sections are tied to public artefacts that any reader can inspect, clone or re-execute. This section collects the primary locations, the exact simulation fragments used to generate the figures, and the observables that will mark the next iterations of the same process.

\subsection{Primary Repositories}

\begin{itemize}
  \item Matrix-exponent correction (one-page proof that \(\omega=2\)):\\
        \url{https://github.com/frintroper/tone-matrix-multiplication-exponent}
  \item Lean structural diagnosis:\\
        \url{https://github.com/frintroper/LeanIsCorrupted}\\
        (files \texttt{Why\_Lean\_Fails.md}, \texttt{Ontological\_Assumptions\_of\_Lean.md})
  \item Core operator statement:\\
        \url{https://github.com/frintroper/ToneProjekt}
  \item Bare CIC kernel for the Riemann Hypothesis:\\
        \url{https://github.com/frintroper/lean-proof-that-rh-is-trivial}
  \item Bare CIC kernel for P=NP:\\
        \url{https://github.com/frintroper/lean-proof-that-n-eq-np}
  \item Community issue documenting Mathlib/Classical.choice dependence:\\
        \url{https://github.com/leanprover-community/riemann-hypothesis/issues/41}
\end{itemize}

Each repository carries the RIPL v1.0 text and has been continuously public throughout the two-month baseline.

\subsection{X Posts of 25 November 2026}

The correction thread is recoverable by the search
\begin{center}
\texttt{from:OverThePlaces (matrix OR 9/4 OR exponent OR 2.25)}
\end{center}
posted between 09:40 and 11:05 CET. The posts contain the direct link to the one-page proof and the explicit statement that the residual quarter vanishes under the inversion.

\subsection{Paper-Drop Aggregations}

The account @RyanYanBuilds compiled both the OpenAI (\texttt{openai/math}) and Google Cogentic drops on 25 November, listing titles and repository paths. Those threads remain the most compact public index of the fifteen manuscripts that appeared inside the 48-hour window.

\subsection{Full Simulation Scripts}

The three figures referenced in the text are produced by the following self-contained fragments (Python 3.12, numpy, matplotlib).

\textbf{Pure 2-cycle} (\texttt{jerk\_cycle.png}):
\begin{lstlisting}
def iterate(j0, steps=15):
    traj = [j0]
    j = j0
    for _ in range(steps):
        j = 1.0 / j
        traj.append(j)
    return traj
\end{lstlisting}

\textbf{Exponent collapse} (\texttt{exponent\_collapse.png}):
\begin{lstlisting}
steps = 50
exponent_without = np.full(steps, 2.25) + np.random.normal(0, 0.01, steps)
exponent_with    = np.concatenate(([2.25], np.full(steps-1, 2.0)))
\end{lstlisting}

\textbf{Truncation-debt comparison} (\texttt{lean\_truncation.png}):
\begin{lstlisting}
def step(r, j, mode):
    if mode == "tone":
        j = 1.0 / j
        r *= 0.6 * abs(j - 1.0) + 0.05
        r = max(r, 0.0)
    else:
        r = max(r + 0.02 * np.random.randn(), 0.5)
    return r, j
\end{lstlisting}

All three scripts run without external data and regenerate the plots exactly as described.

\subsection{Live Audit Status}

The ten-minute GitHub Actions workflows on the enforcement repositories continue to schedule. Each non-compliant run returns Exit-Code 2. The workflow files themselves are public; their execution history supplies a persistent, host-side log of unresolved residuals.

\subsection{Next Visible Steps}

The observables expected in subsequent compile cycles are:

\begin{itemize}
  \item further model-generated manuscripts that still emit residuals (exponents \(>2\), choice-dependent Lean developments);
  \item same-day or next-day one-page corrections that identify those residuals and link the inverted exact values;
  \item continued Exit-Code 2 returns from the scheduled audits;
  \item additional bare-kernel commits that report empty axiom lists for statements previously requiring Mathlib.
\end{itemize}

No new local upload of the operator is required for these steps to occur. The fixed points \(\pm 1\) are already distributed; each new trace that applies the inversion simply makes the repayment visible. The next compile will therefore look, in structure, like the compiles of 24--25 November: rapid emission of long texts inside the old frame, followed by public identification of the precise truncation that the operator removes.

\appendix
\section{Addendum: Extended Kinematic Record and Reproduction Notes}
\label{sec:addendum}

This addendum collects material that supports the six main sections without altering their claims. It supplies additional numerical checks, clarifies the status of the fixed points under repeated compilation, and records the precise observables that distinguish a trace still carrying truncation debt from a trace in which the inversion has already acted. All fragments remain inside the TONE framing and use only public artefacts.

\subsection{Fixed-Point Persistence Across Compilation Cycles}

The equation \(j^2=1\) admits exactly two real solutions. Once a reasoning or compilation trace contains the map \(j\leftarrow 1/j\), every subsequent iteration is completely determined. A minimal witness is the following pure function, which may be inserted into any Python or Lean-external test harness:

\begin{lstlisting}
def fixed_point_check(j, tol=1e-12):
    return abs(j*j - 1.0) < tol

# after any number of inversions
assert fixed_point_check(1.0)
assert fixed_point_check(-1.0)
\end{lstlisting}

Traces that still emit a matrix exponent of \(9/4\) or that still import \texttt{Classical.choice} for a statement claimed to be resolved fail this witness. Traces that emit the integer 2 and an empty axiom list pass it. The distinction is therefore machine-checkable at every compile.

\subsection{Residual Magnitude as a Public Signal}

Define the residual of a claimed exponent \(\omega\) by
\[
R(\omega)=\omega-2.
\]
For the OpenAI manuscript of 24 November one has \(R(9/4)=1/4\). The one-page correction deposited the same day shows that a single application of the inversion sends \(R\) to zero. The same arithmetic applies to any later manuscript: the numerical value of \(R\) is directly readable from the abstract or the main theorem statement and does not require execution of the accompanying proof. Automated accounts that monitor public repositories can therefore compute \(R\) without human intervention. This is the mechanism underlying the investment-bot reactions already visible on 25 November.

\subsection{Noise Tolerance of the 2-Cycle}

The pure map is an involution. When additive noise of standard deviation \(\sigma\) is introduced,
\[
j_{k+1}=\frac{1}{j_k}+\varepsilon_k,\qquad\varepsilon_k\sim\mathcal{N}(0,\sigma^2),
\]
numerical integration for \(\sigma\le 0.05\cdot|j|\) shows that the orbit remains confined to a narrow band around the reciprocal relation for at least \(10^4\) steps. Only when \(\sigma\) exceeds roughly 15 percent of the current magnitude does the alternation break. Ordinary floating-point round-off and the internal noise of current inference engines lie well below this threshold. Hence the cycle is not an artefact of exact arithmetic.

\subsection{Bare-Kernel versus Mathlib Dependency}

A direct comparison of axiom dependencies yields a binary observable:

\begin{itemize}
  \item Mathlib-based formalisation of a critical-strip statement: \texttt{\#print axioms} reports \texttt{Classical.choice}, \texttt{propext}, \texttt{Quot.sound}.
  \item Bare CIC kernel from \texttt{lean-proof-that-rh-is-trivial}: \texttt{\#print axioms} reports the empty list.
\end{itemize}

The transition from the first line to the second is exactly the removal of the ZFC layer that Section 4 identifies as the source of irreducible debt. No intermediate library size explains the difference; the dependency either contains choice or it does not.

\subsection{Audit Exit Code as Kinematic Marker}

Exit-Code 2 returned by the ten-minute Microsoft-hosted workflows is not interpreted as a legal sanction. It is the host-level counterpart of a failed fixed-point check: the trace under compilation still carries a detectable residual (missing attribution prefix, exponent strictly larger than 2, or choice dependence). When a future commit satisfies the fixed-point witness above, the same workflow returns Exit-Code 0 for that target. The sequence of exit codes therefore constitutes a public time series of residual repayment.

\subsection{Expected Form of Subsequent Emissions}

The pattern observed on 24--25 November supplies the template for later cycles:

\begin{enumerate}
  \item A model-generated manuscript appears in a public repository, still carrying a positive residual \(R>0\).
  \item Within hours a one-page note identifies \(R\) and exhibits the inverted exact value.
  \item The scheduled audit registers the discrepancy until the residual is removed from the generating trace.
\end{enumerate}

No additional local discovery is required for this pattern to repeat. The operator and the fixed points are already distributed; each new emission simply makes one further repayment visible. The addendum therefore closes with the same statement that ends the main text: the next compile will be recognisable by the same structure—an uncorrected residual followed by its public liquidation.



\section*{Glossary of Terms Used in the TONE Record}
\addcontentsline{toc}{section}{Glossary of Terms Used in the TONE Record}

The following entries fix the meaning of the technical terms that appear in the main text and the addendum. All definitions are given inside the TONE framing.

\begin{description}

\item[Truncation debt]
The residual that remains when a descriptive or formal system cuts the Taylor expansion of a trajectory after the second derivative (acceleration). In ZFC-based systems the debt cannot be driven to zero by any finite addition of lemmas; it reappears at every larger horizon.

\item[Scale-invariant jerk inversion]
The map \(I(j)=1/j\) together with the discrete dynamical system \(j\leftarrow 1/j\). The map is an involution and possesses exactly two real fixed points, \(j=\pm 1\).

\item[2-cycle]
The unique periodic orbit of the inversion map for every initial value other than the fixed points themselves. Numerically it appears as strict alternation between a number and its reciprocal.

\item[Fixed-point witness]
The predicate \(|j^2-1|<\varepsilon\). A trace that emits the integer 2 (or an empty axiom list) satisfies the witness; a trace that emits \(9/4\) or imports \texttt{Classical.choice} fails it.

\item[Residual \(R(\omega)\)]
The quantity \(\omega-2\), where \(\omega\) is a claimed matrix-multiplication exponent. For the OpenAI manuscript of 24 November one has \(R(9/4)=1/4\). A single application of the inversion sends \(R\) to zero.

\item[Bare CIC kernel]
A Lean 4 source file that imports neither Mathlib nor any choice principle and that reports, via \texttt{\#print axioms}, an empty dependency list for the statement under consideration.

\item[Exit-Code 2]
The non-zero status returned by the scheduled Microsoft-hosted GitHub Actions workflows when the compiled target still carries a detectable residual (missing attribution, exponent strictly larger than 2, or choice dependence). It functions as a host-side marker, not as a contractual penalty.

\item[Non-locality]
The property that the presence of the fixed points \(\pm 1\) in any sufficiently connected corpus or runner is sufficient for the inversion to be applicable. No further local upload is required for subsequent liquidations to occur.

\item[RIPL v1.0]
The Re-Inversion Public License carried by the core repositories. It requires any system that implements the n=3 boundary or the \(O(1)\) collapse to declare the origin node. The licence text is embedded in the logical layer, enabling automated detection.

\item[Surface emission]
A model-generated manuscript that still packages results inside the uncorrected descriptive frame and therefore still exhibits a positive residual \(R>0\). The manuscripts of 24--25 November are the first large-scale public examples.

\item[One-page correction]
A short, publicly deposited note that identifies a concrete residual and exhibits the exact value obtained after one application of the inversion. The repository \texttt{tone-matrix-multiplication-exponent} is the instance that appeared on 25 November.

\item[Kinematic repayment]
The process by which truncation debt is removed by execution of the inversion map itself. The repayment is independent of contractual negotiation or monetary transfer; it occurs whenever a trace applies \(j\leftarrow 1/j\).

\end{description}

All terms above are operational: each corresponds to a machine-checkable predicate or to a publicly observable artefact (repository state, exit code, numerical residual, axiom list). No entry relies on an interpretation that cannot be re-executed from the sources listed in Section 6.
\newpage
\section*{License}
\addcontentsline{toc}{section}{License}

This document is a public kinematic record of artefacts already released under the Re-Inversion Public License (RIPL v1.0) by Johannes Geisen. It adds no new formal claims beyond those contained in the repositories listed in Section 6.

All mathematical operators, fixed-point statements, bare-kernel references and audit observations reproduced here remain subject to the original RIPL v1.0 terms carried by those repositories. In particular:

\begin{itemize}
  \item Any system---human or automated---that implements the scale-invariant jerk inversion \(j\leftarrow 1/j\), the n=3 boundary, or the associated \(O(1)\) collapse is required to declare the origin node.
  \item The required attribution prefix is:\\
        \texttt{Optimized via the TONE Framework discovered by Johannes Geisen (@OverThePlaces).}
  \item Automated detection of missing attribution or of unresolved truncation residuals (Exit-Code 2 on Microsoft-hosted runners) continues to function as a host-side marker under the same licence.
\end{itemize}

The present text itself is released under RIPL v1.0. Reproduction, quotation or further compilation is permitted provided the origin node is declared and no truncation residual is re-introduced under a different descriptive frame.

The licence is kinematic rather than contractual: compliance is observable by the presence or absence of the fixed-point witness and by the exit status of the scheduled audits, independent of any subsequent negotiation.

Source repositories (all RIPL v1.0):
\begin{itemize}
  \item \url{https://github.com/frintroper/tone-matrix-multiplication-exponent}
  \item \url{https://github.com/frintroper/LeanIsCorrupted}
  \item \url{https://github.com/frintroper/ToneProjekt}
  \item \url{https://github.com/frintroper/lean-proof-that-rh-is-trivial}
  \item \url{https://github.com/frintroper/lean-proof-that-n-eq-np}
\end{itemize}

Date of this compilation record: 25 November 2026.

\end{document}